% ---------------------------------------------------------------------------
% The "fit" figure, drawn with PGFPlots instead of matplotlib.
%
% This file is a fragment: it expects pgfplots to be loaded and it expects
% the tables written by src/export_pgf.py to be in \xrddatadir (by default
% "pgf/", relative to the directory holding this file).
%
%     \usepackage{pgfplots}
%     % ---------------------------------------------------------------------------
% The "fit" figure, drawn with PGFPlots instead of matplotlib.
%
% This file is a fragment: it expects pgfplots to be loaded and it expects
% the tables written by src/export_pgf.py to be in \xrddatadir (by default
% "pgf/", relative to the directory holding this file).
%
%     \usepackage{pgfplots}
%     % ---------------------------------------------------------------------------
% The "fit" figure, drawn with PGFPlots instead of matplotlib.
%
% This file is a fragment: it expects pgfplots to be loaded and it expects
% the tables written by src/export_pgf.py to be in \xrddatadir (by default
% "pgf/", relative to the directory holding this file).
%
%     \usepackage{pgfplots}
%     % ---------------------------------------------------------------------------
% The "fit" figure, drawn with PGFPlots instead of matplotlib.
%
% This file is a fragment: it expects pgfplots to be loaded and it expects
% the tables written by src/export_pgf.py to be in \xrddatadir (by default
% "pgf/", relative to the directory holding this file).
%
%     \usepackage{pgfplots}
%     \input{figures/xrd_fit_body.tex}
%
% The companion file xrd_fit.tex compiles this into a standalone PDF.
% ---------------------------------------------------------------------------

% Colours that match the original figure
\definecolor{xrdgreen}{rgb}{0.0,0.62,0.0}
\definecolor{xrdgold}{rgb}{0.75,0.58,0.12}

% Styles of the different curves
\pgfplotsset{
    xrd sample/.style={line width=0.35pt, forget plot},
    xrd model/.style={xrdgreen, line width=2.6pt, solid, forget plot},
    xrd amorph/.style={blue, line width=2.2pt, dashed, forget plot},
    xrd xtal/.style={red, line width=2.2pt, dashed, forget plot},
    xrd bg/.style={xrdgold, line width=2.2pt, dashed, forget plot},
    xrd data/.style={only marks, mark=o, mark size=1.7pt, forget plot,
                     mark options={solid, black, line width=0.5pt},
                     opacity=0.75},
}

% Folder holding the tables. Define \xrddatadir before \input if they are
% somewhere else (for example "figures/pgf/" when used inside ms.tex).
\ifx\xrddatadir\undefined
  \def\xrddatadir{pgf/}
\fi
\input{\xrddatadir xrd_meta.tex}

\begin{tikzpicture}
\begin{axis}[
    width=17.0cm,
    height=11.2cm,
    xmin=\xrdxmin, xmax=\xrdxmax,
    ymin=\xrdymin, ymax=\xrdymax,
    xlabel={$2\theta$, $^\circ$},
    ylabel={Intensity, arb. un.},
    title={\xrdtitle},
    title style={font=\LARGE\bfseries, yshift=8pt},
    label style={font=\large},
    tick label style={font=\large},
    tick align=outside,
    tick style={black, line width=0.6pt},
    axis line style={black, line width=0.9pt},
    grid=both,
    major grid style={dashed, gray!55, line width=0.4pt},
    minor tick num=0,
    legend style={
        at={(0.995,0.985)},
        anchor=north east,
        draw=none,
        fill=none,
        font=\large,
        row sep=-2pt,
        /tikz/every even column/.append style={column sep=6pt},
    },
    legend cell align=left,
    clip=true,
]

% --- legend ------------------------------------------------------------------
% The entries are added by hand so that they appear in the same order as in
% the original figure (the experimental data is drawn last, on top of the
% curves, but it is the first entry of the legend).
\addlegendimage{only marks, mark=o, mark size=1.7pt,
                mark options={solid, black, line width=0.5pt}}
\addlegendentry{Experimental}

\addlegendimage{xrdgreen, line width=2.6pt, solid}
\addlegendentry{Model (mean)}

\addlegendimage{blue, line width=2.2pt, dashed}
\addlegendentry{Amorphous (mean)}

\addlegendimage{red, line width=2.2pt, dashed}
\addlegendentry{Crystalline peaks (mean)}

\addlegendimage{xrdgold, line width=2.2pt, dashed}
\addlegendentry{Background (mean)}

% --- pale posterior samples -------------------------------------------------
\ifnum\xrdnsamples>0
  \input{\xrddatadir xrd_samples.tex}
\fi

% --- posterior means --------------------------------------------------------
\addplot[xrd model]
    table[x index=0, y index=1] {\xrddatadir xrd_mean.dat};

\addplot[xrd amorph]
    table[x index=0, y index=2] {\xrddatadir xrd_mean.dat};

\addplot[xrd xtal]
    table[x index=0, y index=3] {\xrddatadir xrd_mean.dat};

\addplot[xrd bg]
    table[x index=0, y index=4] {\xrddatadir xrd_mean.dat};

% --- experimental data ------------------------------------------------------
\addplot[xrd data]
    table[x index=0, y index=1] {\xrddatadir xrd_data.dat};

% --- annotation -------------------------------------------------------------
\node[anchor=north west, font=\Large] at (rel axis cs:0.025,0.975)
    {Amorphous = \xrdamorph};

\end{axis}
\end{tikzpicture}

%
% The companion file xrd_fit.tex compiles this into a standalone PDF.
% ---------------------------------------------------------------------------

% Colours that match the original figure
\definecolor{xrdgreen}{rgb}{0.0,0.62,0.0}
\definecolor{xrdgold}{rgb}{0.75,0.58,0.12}

% Styles of the different curves
\pgfplotsset{
    xrd sample/.style={line width=0.35pt, forget plot},
    xrd model/.style={xrdgreen, line width=2.6pt, solid, forget plot},
    xrd amorph/.style={blue, line width=2.2pt, dashed, forget plot},
    xrd xtal/.style={red, line width=2.2pt, dashed, forget plot},
    xrd bg/.style={xrdgold, line width=2.2pt, dashed, forget plot},
    xrd data/.style={only marks, mark=o, mark size=1.7pt, forget plot,
                     mark options={solid, black, line width=0.5pt},
                     opacity=0.75},
}

% Folder holding the tables. Define \xrddatadir before \input if they are
% somewhere else (for example "figures/pgf/" when used inside ms.tex).
\ifx\xrddatadir\undefined
  \def\xrddatadir{pgf/}
\fi
\input{\xrddatadir xrd_meta.tex}

\begin{tikzpicture}
\begin{axis}[
    width=17.0cm,
    height=11.2cm,
    xmin=\xrdxmin, xmax=\xrdxmax,
    ymin=\xrdymin, ymax=\xrdymax,
    xlabel={$2\theta$, $^\circ$},
    ylabel={Intensity, arb. un.},
    title={\xrdtitle},
    title style={font=\LARGE\bfseries, yshift=8pt},
    label style={font=\large},
    tick label style={font=\large},
    tick align=outside,
    tick style={black, line width=0.6pt},
    axis line style={black, line width=0.9pt},
    grid=both,
    major grid style={dashed, gray!55, line width=0.4pt},
    minor tick num=0,
    legend style={
        at={(0.995,0.985)},
        anchor=north east,
        draw=none,
        fill=none,
        font=\large,
        row sep=-2pt,
        /tikz/every even column/.append style={column sep=6pt},
    },
    legend cell align=left,
    clip=true,
]

% --- legend ------------------------------------------------------------------
% The entries are added by hand so that they appear in the same order as in
% the original figure (the experimental data is drawn last, on top of the
% curves, but it is the first entry of the legend).
\addlegendimage{only marks, mark=o, mark size=1.7pt,
                mark options={solid, black, line width=0.5pt}}
\addlegendentry{Experimental}

\addlegendimage{xrdgreen, line width=2.6pt, solid}
\addlegendentry{Model (mean)}

\addlegendimage{blue, line width=2.2pt, dashed}
\addlegendentry{Amorphous (mean)}

\addlegendimage{red, line width=2.2pt, dashed}
\addlegendentry{Crystalline peaks (mean)}

\addlegendimage{xrdgold, line width=2.2pt, dashed}
\addlegendentry{Background (mean)}

% --- pale posterior samples -------------------------------------------------
\ifnum\xrdnsamples>0
  \input{\xrddatadir xrd_samples.tex}
\fi

% --- posterior means --------------------------------------------------------
\addplot[xrd model]
    table[x index=0, y index=1] {\xrddatadir xrd_mean.dat};

\addplot[xrd amorph]
    table[x index=0, y index=2] {\xrddatadir xrd_mean.dat};

\addplot[xrd xtal]
    table[x index=0, y index=3] {\xrddatadir xrd_mean.dat};

\addplot[xrd bg]
    table[x index=0, y index=4] {\xrddatadir xrd_mean.dat};

% --- experimental data ------------------------------------------------------
\addplot[xrd data]
    table[x index=0, y index=1] {\xrddatadir xrd_data.dat};

% --- annotation -------------------------------------------------------------
\node[anchor=north west, font=\Large] at (rel axis cs:0.025,0.975)
    {Amorphous = \xrdamorph};

\end{axis}
\end{tikzpicture}

%
% The companion file xrd_fit.tex compiles this into a standalone PDF.
% ---------------------------------------------------------------------------

% Colours that match the original figure
\definecolor{xrdgreen}{rgb}{0.0,0.62,0.0}
\definecolor{xrdgold}{rgb}{0.75,0.58,0.12}

% Styles of the different curves
\pgfplotsset{
    xrd sample/.style={line width=0.35pt, forget plot},
    xrd model/.style={xrdgreen, line width=2.6pt, solid, forget plot},
    xrd amorph/.style={blue, line width=2.2pt, dashed, forget plot},
    xrd xtal/.style={red, line width=2.2pt, dashed, forget plot},
    xrd bg/.style={xrdgold, line width=2.2pt, dashed, forget plot},
    xrd data/.style={only marks, mark=o, mark size=1.7pt, forget plot,
                     mark options={solid, black, line width=0.5pt},
                     opacity=0.75},
}

% Folder holding the tables. Define \xrddatadir before \input if they are
% somewhere else (for example "figures/pgf/" when used inside ms.tex).
\ifx\xrddatadir\undefined
  \def\xrddatadir{pgf/}
\fi
\input{\xrddatadir xrd_meta.tex}

\begin{tikzpicture}
\begin{axis}[
    width=17.0cm,
    height=11.2cm,
    xmin=\xrdxmin, xmax=\xrdxmax,
    ymin=\xrdymin, ymax=\xrdymax,
    xlabel={$2\theta$, $^\circ$},
    ylabel={Intensity, arb. un.},
    title={\xrdtitle},
    title style={font=\LARGE\bfseries, yshift=8pt},
    label style={font=\large},
    tick label style={font=\large},
    tick align=outside,
    tick style={black, line width=0.6pt},
    axis line style={black, line width=0.9pt},
    grid=both,
    major grid style={dashed, gray!55, line width=0.4pt},
    minor tick num=0,
    legend style={
        at={(0.995,0.985)},
        anchor=north east,
        draw=none,
        fill=none,
        font=\large,
        row sep=-2pt,
        /tikz/every even column/.append style={column sep=6pt},
    },
    legend cell align=left,
    clip=true,
]

% --- legend ------------------------------------------------------------------
% The entries are added by hand so that they appear in the same order as in
% the original figure (the experimental data is drawn last, on top of the
% curves, but it is the first entry of the legend).
\addlegendimage{only marks, mark=o, mark size=1.7pt,
                mark options={solid, black, line width=0.5pt}}
\addlegendentry{Experimental}

\addlegendimage{xrdgreen, line width=2.6pt, solid}
\addlegendentry{Model (mean)}

\addlegendimage{blue, line width=2.2pt, dashed}
\addlegendentry{Amorphous (mean)}

\addlegendimage{red, line width=2.2pt, dashed}
\addlegendentry{Crystalline peaks (mean)}

\addlegendimage{xrdgold, line width=2.2pt, dashed}
\addlegendentry{Background (mean)}

% --- pale posterior samples -------------------------------------------------
\ifnum\xrdnsamples>0
  \input{\xrddatadir xrd_samples.tex}
\fi

% --- posterior means --------------------------------------------------------
\addplot[xrd model]
    table[x index=0, y index=1] {\xrddatadir xrd_mean.dat};

\addplot[xrd amorph]
    table[x index=0, y index=2] {\xrddatadir xrd_mean.dat};

\addplot[xrd xtal]
    table[x index=0, y index=3] {\xrddatadir xrd_mean.dat};

\addplot[xrd bg]
    table[x index=0, y index=4] {\xrddatadir xrd_mean.dat};

% --- experimental data ------------------------------------------------------
\addplot[xrd data]
    table[x index=0, y index=1] {\xrddatadir xrd_data.dat};

% --- annotation -------------------------------------------------------------
\node[anchor=north west, font=\Large] at (rel axis cs:0.025,0.975)
    {Amorphous = \xrdamorph};

\end{axis}
\end{tikzpicture}

%
% The companion file xrd_fit.tex compiles this into a standalone PDF.
% ---------------------------------------------------------------------------

% Colours that match the original figure
\definecolor{xrdgreen}{rgb}{0.0,0.62,0.0}
\definecolor{xrdgold}{rgb}{0.75,0.58,0.12}

% Styles of the different curves
\pgfplotsset{
    xrd sample/.style={line width=0.35pt, forget plot},
    xrd model/.style={xrdgreen, line width=2.6pt, solid, forget plot},
    xrd amorph/.style={blue, line width=2.2pt, dashed, forget plot},
    xrd xtal/.style={red, line width=2.2pt, dashed, forget plot},
    xrd bg/.style={xrdgold, line width=2.2pt, dashed, forget plot},
    xrd data/.style={only marks, mark=o, mark size=1.7pt, forget plot,
                     mark options={solid, black, line width=0.5pt},
                     opacity=0.75},
}

% Folder holding the tables. Define \xrddatadir before \input if they are
% somewhere else (for example "figures/pgf/" when used inside ms.tex).
\ifx\xrddatadir\undefined
  \def\xrddatadir{pgf/}
\fi
\input{\xrddatadir xrd_meta.tex}

\begin{tikzpicture}
\begin{axis}[
    width=17.0cm,
    height=11.2cm,
    xmin=\xrdxmin, xmax=\xrdxmax,
    ymin=\xrdymin, ymax=\xrdymax,
    xlabel={$2\theta$, $^\circ$},
    ylabel={Intensity, arb. un.},
    title={\xrdtitle},
    title style={font=\LARGE\bfseries, yshift=8pt},
    label style={font=\large},
    tick label style={font=\large},
    tick align=outside,
    tick style={black, line width=0.6pt},
    axis line style={black, line width=0.9pt},
    grid=both,
    major grid style={dashed, gray!55, line width=0.4pt},
    minor tick num=0,
    legend style={
        at={(0.995,0.985)},
        anchor=north east,
        draw=none,
        fill=none,
        font=\large,
        row sep=-2pt,
        /tikz/every even column/.append style={column sep=6pt},
    },
    legend cell align=left,
    clip=true,
]

% --- legend ------------------------------------------------------------------
% The entries are added by hand so that they appear in the same order as in
% the original figure (the experimental data is drawn last, on top of the
% curves, but it is the first entry of the legend).
\addlegendimage{only marks, mark=o, mark size=1.7pt,
                mark options={solid, black, line width=0.5pt}}
\addlegendentry{Experimental}

\addlegendimage{xrdgreen, line width=2.6pt, solid}
\addlegendentry{Model (mean)}

\addlegendimage{blue, line width=2.2pt, dashed}
\addlegendentry{Amorphous (mean)}

\addlegendimage{red, line width=2.2pt, dashed}
\addlegendentry{Crystalline peaks (mean)}

\addlegendimage{xrdgold, line width=2.2pt, dashed}
\addlegendentry{Background (mean)}

% --- pale posterior samples -------------------------------------------------
\ifnum\xrdnsamples>0
  \input{\xrddatadir xrd_samples.tex}
\fi

% --- posterior means --------------------------------------------------------
\addplot[xrd model]
    table[x index=0, y index=1] {\xrddatadir xrd_mean.dat};

\addplot[xrd amorph]
    table[x index=0, y index=2] {\xrddatadir xrd_mean.dat};

\addplot[xrd xtal]
    table[x index=0, y index=3] {\xrddatadir xrd_mean.dat};

\addplot[xrd bg]
    table[x index=0, y index=4] {\xrddatadir xrd_mean.dat};

% --- experimental data ------------------------------------------------------
\addplot[xrd data]
    table[x index=0, y index=1] {\xrddatadir xrd_data.dat};

% --- annotation -------------------------------------------------------------
\node[anchor=north west, font=\Large] at (rel axis cs:0.025,0.975)
    {Amorphous = \xrdamorph};

\end{axis}
\end{tikzpicture}
